\chapter{Trascrizione utilizzatori}\label{app:risultati-interviste-utilizzatori}
In questa appendice sono riportate le trascrizioni delle interviste condotte con gli utilizzatori abituali
del sistema attuale, al fine di raccogliere i requisiti per il nuovo software.

\section{Operazioni comuni}
Le operazioni più comuni svolte all'interno dell'applicativo includono:
\begin{itemize}
      \item Creazione di offerte per ricambi e assistenza (\ac{RA}).
      \item Creazione di offerte per clienti non diretti.
      \item Creazione delle offerte Access per conto di agenti o rivenditori esterni.
      \item Creazione di offerte per clienti diretti.
      \item Creazione degli ordini ogni qualvolta un'offerta viene accettata.
      \item Ricerca e copia di vecchie offerte simili durante la creazione di nuove offerte.
      \item Correzione delle traduzioni preimpostate fornite da Access.
\end{itemize}

\section{Gestione protezione dati e accessi}
\subsection{Gestione utenti e visibilità degli agenti}
Per quanto riguarda la gestione degli utenti e la visibilità dei dati, gli agenti
hanno evidenziato i seguenti punti chiave:
\begin{itemize}
      \item \textbf{Segregazione dei dati:} È fondamentale mantenere la distinzione tra offerte e
            \ac{RA} (Richieste assistenza/ricambi).
            La distinzione tra i due tipi aiuta a evitare confusione e garantisce che le operazioni
            siano gestite correttamente.
      \item \textbf{Privacy per autore:} Un agente esterno non deve vedere i progetti o le offerte
            degli altri agenti per evitare conflitti di interesse (Policy aziendale).
      \item \textbf{Anagrafica clienti condivisa:} A differenza delle offerte, l'elenco clienti deve
            essere visibile a tutti, può capitare di fare offerte per conto di altri commerciali.
\end{itemize}

\section{Flussi di lavoro attuali}

\subsection{Ricambi e assistenza post-vendita}
Il flusso di lavoro attuale, per la gestione dei ricambi e dell'assistenza
post-vendita, presenta diverse inefficienze dovute alla frammentazione degli strumenti
e alla mancanza di dati strutturati.

\begin{itemize}
      \item \textbf{Template di partenza:} Attualmente non esiste un template dedicato; si parte quindi
            da vecchi ordini e, per individuare quali siano le assistenze, l'unica discriminante è
            l'autore.
            Questa pratica permette di ereditare tutte le condizioni di vendita e i dati del cliente,
            risparmiando tempo nella compilazione (operazione attualmente molto onerosa in termini di tempo).
      \item \textbf{Ricerca ricambi manuale:} L'operatore deve consultare disegni tecnici (Fusion, PDF, cartacei)
            e database separati (Access, Ad-Hoc, Excel, Server ``Commesse'') per identificare i codici dei ricambi.
      \item \textbf{Inserimento del componente:} Una volta individuato il componente, non essendoci una
            corrispondenza univoca dei codici, l'operatore deve trovare la corrispondenza in
            Access o creare un nuovo codice per il ricambio.
      \item \textbf{Creazione offerta complessa:} Non è facile né la ricerca né l'assegnazione di un prezzo
            se il componente è datato; la compilazione dell'offerta talvolta richiede calcoli manuali (prezzi, pesi, volumi).
            Spesso si ricorre alla copia di vecchie offerte \ac{RA} o commesse, pratica che trascina errori pregressi.
\end{itemize}

\subsection{Impianti nuovi e ampliamenti}
L'attuale processo di creazione delle offerte prevede i seguenti passaggi:
\begin{enumerate}
      \item \textbf{Raccolta requisiti:} Il commerciale si interfaccia con il cliente per comprendere le
            esigenze tecniche e commerciali.
      \item \textbf{Ricerca e copiatura di una vecchia offerta:} Spesso si parte da un'offerta precedente simile
            per velocizzare il processo.
      \item \textbf{Selezione prodotti:} Aggiunta manuale dei prodotti, uno a uno, tramite descrizione.
      \item \textbf{Configurazione offerta:} Creazione dell'offerta suddivisa in capitoli, applicando gli sconti
            e i calcoli necessari.
      \item \textbf{Invio al cliente:} Generazione del documento finale in formato PDF o cartaceo e invio al cliente.
      \item \textbf{Approvazione offerta:} In caso di accettazione, si procede con la conferma d'ordine e
            la comunicazione all'ufficio produzione e all'ufficio amministrativo.
\end{enumerate}

\section{Problemi attuali del sistema}
Dall'intervista i commerciali hanno fatto emergere le seguenti criticità del sistema attuale:
\begin{itemize}
      \item \textbf{Lentezza e complessità:} L'intero processo di gestione è lento e laborioso
            e richiede molto tempo per completare un'offerta. Questo è dovuto alla necessità di consultare molteplici fonti di
            dati non integrate. Il programma stesso utilizzato (Access) risulta lento e macchinoso, con molti tempi morti.
      \item \textbf{Disallineamento codici:} Esiste una discrepanza tra i codici commerciali e quelli tecnici/gestionali.
            Questo costringe a ricerche manuali su schemi funzionali o database separati per identificare il pezzo corretto (Ad-Hoc, Fusion, \dots).
      \item \textbf{Mancanza di automazione nei calcoli:}
            \begin{itemize}
                  \item \textbf{Prezzi:} I costi e i prezzi di vendita devono essere spesso calcolati manualmente o verificati
                        su Excel/Ad-Hoc, specialmente per componenti su misura o obsoleti.
                  \item \textbf{Logistica:} Non esiste un calcolo automatico di volumi e pesi per il trasporto; l'utente deve
                        ipotizzare i dati basandosi su vecchi DDT o sull'esperienza. Il costo è talvolta calcolato mediante il sito del corriere.
                  \item \textbf{Manodopera:} I tempi e i costi per gli interventi tecnici (es. sostituzione saccone) non sono
                        standardizzati e richiedono consultazioni continue con i colleghi (es. Logistica).
            \end{itemize}
      \item \textbf{Componenti elettrici:} La gestione e l'inserimento di materiale elettrico sono molto difficili per la 
      complessità e la varietà di codici esistenti. Attualmente si consulta direttamente l'ufficio tecnico elettrico per
            avere un elenco dei componenti necessari da offrire.
      \item \textbf{Gestione obsolescenza:} È difficile sapere se un componente venduto anni fa è stato sostituito da una
            nuova versione o se è obsoleto, rischiando di vendere pezzi non più procurabili. Serve un sistema di tracciamento delle
            sostituzioni e delle obsolescenze (es. Fusion).
      \item \textbf{Discrepanza ``As Built'':} Le modifiche fatte in cantiere (\ac{RA} di cantiere) o successive alla vendita non
            sono sempre tracciate in modo unitario, rendendo difficile conoscere la configurazione reale dell'impianto attuale, punto
            fondamentale per una corretta assistenza senza errori.
      \item \textbf{Sistema di traduzione:} È fondamentale avere un sistema di traduzione
            integrato e affidabile per le descrizioni
            dei componenti e delle offerte; attualmente, oltre alle traduzioni preimpostate di Access
            (spesso da correggere), si utilizzano
            servizi esterni (Google Translate) che richiedono tempo e possono introdurre errori.
\end{itemize}

\section{Requisiti Funzionali per il nuovo software}

\subsection{Database e integrazione}
Per migliorare l'integrazione dei dati e la gestione delle informazioni, il nuovo sistema deve prevedere:
\begin{itemize}
      \item \textbf{Database unificato:} Il nuovo sistema deve integrare o dialogare con Ad-Hoc \ac{ERP}, Fusion \ac{PLM}
            e l'archivio storico per garantire l'univocità del dato.
      \item \textbf{Codifica condivisa:} Allineamento tra codici commerciali e tecnici.
            Il codice usato per l'offerta deve permettere di risalire automaticamente a costi, distinta base e disegni.
      \item \textbf{Gestione revisioni:} Il sistema deve tracciare le revisioni dei clienti
            (cambi nome, stabilimento) e dei componenti (sostituzioni, obsolescenza).
      \item \textbf{Codici a catalogo:} Il sistema deve già contenere tutti i codici esistenti;
            la volontà del commerciale è di non doverli inserire manualmente. È necessario prevedere una
            gestione esterna dedicata all'aggiunta e all'aggiornamento di nuovi codici.
      \item \textbf{Integrazione nel flusso dell'ordine:} Attualmente esistono 3 codici fondamentali per
            la gestione dell'ordine:
            \begin{itemize}
                  \item \textbf{Codice offerta:} Codice univoco dell'offerta commerciale.
                  \item \textbf{Codice ordine:} Codice univoco dell'offerta trasformata in ordine.
                  \item \textbf{Codice distinta:} Codice univoco della distinta base gestionale dell'ordine.
            \end{itemize}
            Il nuovo sistema deve collegare e mostrare questi 3 codici in modo chiaro e visibile per l'operatore,
            facilitando le ricerche e la tracciabilità dell'ordine.
\end{itemize}

\subsection{Ricerca e navigazione ricambi}
Per facilitare la ricerca e la selezione dei ricambi, il sistema deve prevedere:
\begin{itemize}
      \item \textbf{Ricerca multicriterio:} Necessità di cercare per Numero di Commessa (fondamentale),
            Cliente, Stabilimento e Tag Tecnici (es. ``valvola a farfalla'').
      \item \textbf{Struttura ad albero/gerarchica:} Possibilità di esplorare la distinta base \ac{BOM}
            scendendo nei livelli (esplosione componenti) per individuare il singolo ricambio.
      \item \textbf{Filtri per ``Ricambi tipici'':} Implementazione di filtri che mostrino, per una data
            famiglia di macchine (es. Lift, Vibrovaglio), solo i componenti tipicamente venduti come ricambi;
            i codici devono essere allineati ai codici tecnici, velocizzando l'assistenza una volta individuato il componente.
\end{itemize}

\subsection{Creazione e gestione offerte}
Per semplificare e velocizzare la creazione delle offerte di ricambi e assistenza, il sistema deve prevedere:
\begin{itemize}
      \item \textbf{Automazione descrizioni e lingue:} Compilazione automatica delle descrizioni in lingua
            (traduzioni pre-caricate) senza dover usare traduttori esterni.
      \item \textbf{Calcolo prezzi automatizzato:} Il sistema deve recuperare il costo industriale da Ad-Hoc e
            applicare le maggiorazioni o i listini specifici per cliente/settore.
      \item \textbf{Tracciabilità univoca della pratica:} Mantenimento di un riferimento univoco che segua
            la pratica dalla richiesta iniziale all'ordine confermato, per evitare duplicazioni e facilitare la
            rintracciabilità, affiancato da uno stato di avanzamento \textit{(es. ``In Revisione'', ``Inviata'', ``Accettata'',
                  ``Rifiutata'')} per facilitarne l'interpretazione.
      \item \textbf{Distinzione tipologie RA:} Gestione differenziata ma tracciabile per ``Ricambi'',
            ``Ampliamenti'' e ``RA di Cantiere''. Al momento le assistenze sopra citate sono raggruppate sotto
            l'unico cappello del ``VSS''.
\end{itemize}

\subsection{Funzionalità avanzate e UX}
Per migliorare l'esperienza utente e l'efficienza operativa, il sistema deve prevedere:
\begin{itemize}
      \item \textbf{Intestazione offerta e dati cliente:} È importante riportare tutte le informazioni
            attuali del cliente (Indirizzo, Nome, Numero di commessa associata) in modo chiaro e visibile nella
            schermata dell'offerta.
      \item \textbf{Schermata dedicata allo storico:} Una sezione dove poter visualizzare tutte le offerte
            e gli ordini precedenti associati al cliente o alla commessa selezionata.
      \item \textbf{Schermata offerta intuitiva:} Interfaccia user-friendly che guidi l'utente passo-passo
            nella compilazione dell'offerta, minimizzando errori e omissioni.
      \item \textbf{Visualizzazione dati completa:} Interfaccia che mostri tutte le informazioni rilevanti
            (es. dettagli tecnici completi come ``filtrante/non filtrante'') senza dover scorrere o aprire popup.
      \item \textbf{Modalità offline/sync:} Il software deve permettere di lavorare in locale (offline) durante le visite dai clienti
            e sincronizzarsi automaticamente con il database centrale una volta ristabilita la connessione VPN.
      \item \textbf{Multi-finestra:} Possibilità di tenere aperte più schede o finestre di offerte diverse contemporaneamente.
            Utile la possibilità di un sistema multi-documento \textit{Es: AutoCAD}.
      \item \textbf{Interfaccia mobile/tablet:} Ottimizzazione per l'uso su tablet (iPad) per la compilazione
            di dati preliminari e note durante il sopralluogo o in fiera.
\end{itemize}

\section{Note di implementazione}
Gli utilizzatori hanno fornito alcune indicazioni utili per la fase di implementazione:
\begin{itemize}
      \item \textbf{Campi obbligatori:} Rendere obbligatoria la distinzione tra ``Ampliamento'' e ``Ricambi''
            (attualmente campo discrezionale spesso vuoto)
            per migliorare la ricercabilità futura.
      \item \textbf{Beta tester:} Si consiglia di iniziare la fase di test con figure
            tecniche/analitiche e commerciali esperti.
      \item \textbf{Transizione:} Prevista una fase transitoria con utilizzo parallelo dei due sistemi;
            necessaria migrazione o accessibilità facile ai dati storici.
\end{itemize}

\section{Numeri utili}
Per capire gli attuali tempi di lavoro e il volume storico di offerte gestite, si riportano i seguenti dati:
\begin{itemize}
      \item \textbf{Offerte VS Ordini:} Attualmente a sistema risultano create (o presenti nello storico \ac{RA})
            circa 57.000 offerte (\ac{RA}), di cui circa 41.000 si sono trasformate in ordini effettivi.
      \item \textbf{Tempo base per una semplice offerta:} Attualmente, per una richiesta tipo di cui si conoscono già
            il numero di commessa e i codici da inserire, il tempo medio per la creazione di un'offerta è di circa 15 minuti.
            Il tempo sale notevolmente se devono essere ricercate le informazioni mancanti.
\end{itemize}